% 1-implementare.tex starts here
\chapter{Implementare}\label{ch:impl}

\section{Prezentare generală}
Aplicația \textit{Kookio} a fost dezvoltată conform arhitecturii \textbf{PAN} descrise în 
\textbf{Anexa~\ref{ch:appendix-manifest}}. Aceasta integrează tehnologii moderne din ecosistemul
\texttt{TypeScript} pentru a asigura un flux end-to-end unitar între frontend și backend.
În cele ce urmează sunt prezentate componentele interne ale sistemului, fluxurile principale și
mecanismele de automatizare.

\section{Arhitectura sistemului}
\subsection{Structura monorepo-ului Nx}\label{subsec:monorepo-struct}
\vspace{0.3cm}
\begin{TreeVerb}
kookio/
├── web/                        # aplicația Analog
│   │                             (frontend + server SSR)
│   └── src/
│       ├── app/                # componente Angular, pagini, servicii
│       └── server/             # rute Nitro, tRPC și handler-e API
├── libs/
│   ├── data-access/prisma/     # schema Prisma, migrare,
│   │                             client, validări Zod
│   └── shared/                 # mock-uri și utilitare comune
└── nx.json, package.json, etc. # configurarea monorepo-ului Nx
\end{TreeVerb}
\vspace{0.3cm}
\subsection{Flux aplicație end-to-end}\label{subsec:e2e-flow}
Fluxul aplicației urmează schema: \textit{Analog (UI/SSR)} → \textit{tRPC (API type-safe)} → \textit{Prisma (bază de date)} → \textit{Nx (orchestrare)}.
Acest flux permite un lanț complet de tipuri, logare unitară și construcție incrementală a proiectului.
% 1-implementare.tex ends here